% Work draft after F.3.24, for independent audit; full LM6 A.22 unchanged.
\section*{Actual rational CM orbits and the universal torus projector}

\subsection*{The projector after the actual simultaneous constraints}
Keep the genus-three universal submodule $U_6\subset\mathfrak h_6$
and quotient $q_6$ of F.3.24. Fix any symplectic basis of $V$.
Let $T_{\mathrm{all}}$ be the full subspace of special degree-six
derivations tangent to all six basis lines, without imposing a
character condition. Then
\[
                         q_6(U_6\cap T_{\mathrm{all}})=0.       \tag{QCM.1}
\]

Indeed, let $T_{\mathrm{diag}}$ be the full split diagonal
similitude torus for this basis, with multiplier $t$. The
subspace $T_{\mathrm{all}}$ is stable under it: for a diagonal
element $h$, conjugation sends $D(v)=[A_v,v]$ to another
bracket with the same line $\mathbf Q_\ell v$. The universal
submodule $U_6$ is stable under $\operatorname{GSp}(V)$.
Its symplectic stability was proved in F.3.24; scalar matrices
act on degree-six derivations by their sixth power, so this
stability extends across the central isogeny
$\operatorname{Sp}(V)\times\mathbf G_m\to\operatorname{GSp}(V)$.

There is an explicit $\mathbf Q_\ell$-linear weight projector
\[
 e_{t^3}:U_6\cap T_{\mathrm{all}}
       \longrightarrow U_6\cap T_{\mathrm{all}}[t^3].          \tag{QCM.2}
\]
To construct it, choose $h\in T_{\mathrm{diag}}(\mathbf Q_\ell)$
at which $t^3$ differs from every other weight in the finite
dimensional derivation space. Such an $h$ exists because a
finite union of the proper zero sets of the Laurent
polynomials $\lambda-t^3$ does not cover the split torus
over the infinite field $\mathbf Q_\ell$. The Lagrange
interpolation polynomial in the operator $h$ gives (QCM.2).
Every factor preserves both displayed subspaces. The
equivariance of $q_6$ gives
\[
                              q_6 e_{t^3}=q_6.
\]
F.3.24 kills its target and proves (QCM.1).
In an arithmetic application the six tangency constraints
are first proved for an actual Galois element. The operator
(QCM.2) is then applied to that already established linear
constraint space inside the universal image. No lift of
this later algebraic projection to an arithmetic element
is asserted or needed. The same argument applies to the
four-line genus-two quotient calculation in F.3.22.

\subsection*{The rational CM-type identity}
Let $A/\mathbf Q$ be an absolutely simple abelian threefold
with $E=\operatorname{End}^0_{\overline{\mathbf Q}}(A)$ a CM
field of degree six. Write $\Phi$ for the three embeddings
of $E$ appearing in
$\operatorname{Fil}^1H^1_{\mathrm{dR}}(A/\mathbf Q)
 \otimes\overline{\mathbf Q}$.
For $\sigma\in G_{\mathbf Q}$ let
$a_\sigma\in\operatorname{Aut}_{\mathbf Q}(E)$ be its actual
conjugation action on geometric endomorphisms. Then
\[
                            \sigma\Phi=\Phi\circ a_\sigma.    \tag{QCM.3}
\]

Here is the exact semilinear calculation. Choose a field
defining all endomorphisms and their eigenvalues. In the
fixed rational de Rham vector space, let $w$ be an
$E$-eigenvector with character $\varphi\in\Phi$. Applying
$\sigma$ to $e^*w=\varphi(e)w$ gives
\[
             a_\sigma(e)^*\sigma(w)
                        =\sigma(\varphi(e))\sigma(w).
\]
Consequently the character of $\sigma(w)$ is
$\sigma\circ\varphi\circ a_\sigma^{-1}$. The Hodge
filtration is defined over $\mathbf Q$, so it is preserved
by this semilinear map. This proves (QCM.3), including the
direction of the endomorphism automorphism. It is not an
assertion that the arithmetic action on the Tate module
is the abstract permutation action on field embeddings.

\subsection*{Primitive degree-six CM types and the actual endomorphism action}
We first give a direct degree-six type calculation.
Let $E^c$ be the normal closure, $G=\operatorname{Gal}(E^c/\mathbf Q)$,
and choose one embedding in each of the three conjugate
pairs. A CM type is a vector
$\varepsilon\in\{1,-1\}^3$, and $G$ acts on the
three-dimensional anti-invariant space by signed
permutation matrices. Let $W$ be the rational span of
its orbit. Then $W$ has dimension three for the type of
the absolutely simple $A$.

If $W$ had dimension two, its invariant orthogonal
complement would be a rational line. A generator of that
line has three equal nonzero absolute coordinate values:
the finite signed-permutation action preserves length
and is transitive on the three conjugate pairs. Its dot
product with $\varepsilon$ would therefore be a nonzero
multiple of a sum of three numbers in $\{1,-1\}$.
This contradicts orthogonality.

If $W$ had dimension one, write
$g\varepsilon=\eta(g)\varepsilon$ with
$\eta:G\to\{1,-1\}$. Complex conjugation has
$\eta(c)=-1$. The stabilizer of any fixed embedding of
$E$ lies in $\ker\eta$, since that embedding's coordinate
in $\varepsilon$ is nonzero. Hence the imaginary quadratic
field
\[
                    K=(E^c)^{\ker\eta}
\]
is a subfield of $E$. The type $\Phi$ consists exactly of
the embeddings restricting to one fixed embedding of
$K$: the sign $\eta(g)$ determines that restriction.
Thus its Hodge structure is induced from that quadratic
field. This contradicts absolute simplicity by an actual
endomorphism construction. Rational homology is a
three-dimensional $K$-vector space on which the complex
structure is the same scalar complex structure on every
$K$-coordinate. Any $K$-linear projection onto one
$K$-basis coordinate is therefore a rational Hodge
endomorphism. Multiplying it by an integer preserving the
integral lattice gives a complex-linear homomorphism of
the analytic torus $A(\mathbf C)$; its graph is algebraic,
so it is an actual abelian-variety endomorphism. Its image
is an elliptic subvariety, contradicting the absolute
simplicity of $A$. This also gives the exact reason that
the induced quadratic type is unavailable here.

The group $\operatorname{Aut}_{\mathbf Q}(E)$ has order
two or six. Restriction to the totally real cubic
subfield $E^+$ has kernel $\{1,c\}$, and a cubic field
has either one or three automorphisms. If the order
were two, (QCM.3) would place the entire orbit of
$\varepsilon$ in $\{\varepsilon,-\varepsilon\}$,
contradicting the preceding calculation. Thus $E$ is
Galois of degree six. Its complex conjugation is central,
so its Galois group is cyclic, rather than the
nonabelian group of order six.

Let $J$ be the actual image of
$G_{\mathbf Q}\to\operatorname{Aut}_{\mathbf Q}(E)$.
Its order divides six. It has even order: (QCM.3) for
complex conjugation sends $\Phi$ to its complement,+and an odd-order permutation cannot send a set to its
complement. An image of order two would again force
the type orbit to span at most one dimension.
Consequently
\[
 E/\mathbf Q\text{ is cyclic of degree six},\qquad
 G_{\mathbf Q}\twoheadrightarrow\operatorname{Aut}(E)=C_6.
                                                               \tag{QCM.4}
\]

For clarity, the same signed-orbit calculation also
rules out a proper arithmetic character torus in this
simple degree-six CM case. At a completely split prime,
write the six characters as $\chi_\varphi$. Polarization
gives exactly
$\chi_\varphi\chi_{\bar\varphi}=\chi_\ell$.
The proper-smooth comparison B.2.2 and its associated
Hodge--Tate grading apply at every completion without a
reduction assumption. At every embedding of a field
defining all endomorphisms, the exponent-one characters
are the corresponding conjugate of $\Phi$.
A finite-image monomial relation therefore has zero
weight against every conjugate type. Subtracting the
equations for a type and its conjugate, and using the
rank-three signed span, makes its two exponents on
each conjugate pair equal. Adding those equations makes
the sum of the three pair exponents zero. These are
exactly the relations obtained by dividing the three
displayed polarization identities. This character
calculation is additional information; the projector
(QCM.2) does not depend on it.

\subsection*{Actual Tate-module orbit and genus-three recognition}
At any prime where a rational Galois-stable line occurs
after a finite extension, Faltings' theorem identifies
its projector with an idempotent of $E\otimes\mathbf Q_\ell$.
Since $V_\ell A$ is rank one over that algebra, an active
factor must be $\mathbf Q_\ell$. By (QCM.4), $E/\mathbf Q$
is Galois, so all local factors have the same degree.
They therefore all have degree one: the prime splits
completely in $E$. The six eigenspaces are distinct
one-dimensional characters on every sufficiently small
open arithmetic subgroup, since otherwise the
commuting endomorphism algebra would be larger than
$E\otimes\mathbf Q_\ell$.

The actual action on these lines has the exact formula
\[
 \rho_\ell(\sigma)(V_\varphi)
                     =V_{\varphi\circ a_\sigma^{-1}}.       \tag{QCM.5}
\]
Indeed $\rho_\ell(\sigma)$ is $\mathbf Q_\ell$-linear
and conjugates the operator $e$ to $a_\sigma(e)$.
By (QCM.4), the orbit of any one of the six lines is
all six. Rosati duality pairs precisely the conjugate
embedding lines; hence these six lines admit a
symplectic basis, after independently scaling its
members. Formula (QCM.5), rather than abstract
transitivity of embeddings, supplies every orbit line.

It follows that for any smooth proper genus-three
curve $C/\mathbf Q$ with absolutely simple Jacobian,
any rational point $y\in C(\mathbf Q)$, every prime
$\ell$, and every finite number-field extension, a
nontrivial procyclic subgroup of the maximal pro-$\ell$
geometric group of $C\setminus\{y\}$ with open
arithmetic normalizer lies in actual cusp inertia.

Here are the actual maps in the only remaining branch
of the proof. The finite detecting cover and weight
comparison of F.3.21 reduce the obstruction to a
degree-one normalizer line in $V_\ell J_C$. Apply
F.3.18 after defining all endomorphisms. Since $J_C$
is absolutely simple, the nonzero simple quotient
there is isogenous to $J_C$ itself. Its projector-rank
formula gives the CM field $E$ of degree six; thus
(QCM.4)--(QCM.5) apply. Choose six actual arithmetic
lifts giving all six conjugate normalizer subgroups.
Intersect their open images and take a finite Galois
normal core, exactly as in F.3.22--F.3.23. They all
surject onto the resulting common $G_L$.

Nakamura's nonzero universal degree-six quotient and
finite-base-change invariance supply an \emph{actual}
$\sigma\in G_L(6)$ with $q_6(D_\sigma)\ne0$.
Choose a normalizing lift of this same $\sigma$ in
each of the six normalizers. The nonzero degree-one
terms force all six exponents to be one. Removing
the degree-one rational conjugators using the full
logarithms of the subgroup generators, as in F.3.22,
gives $D_\sigma(v)\in[\mathbb L_6(V),v]$ on all six
basis lines. Thus
$D_\sigma\in U_6\cap T_{\mathrm{all}}$.
Equation (QCM.1) contradicts its nonzero quotient.
The degree-one case is excluded. The exact finite-cover
boundary argument and full-normalizer transfer in
F.3.21 now give the cyclotomic character; Nakamura1994,
Theorem2.1.1 then gives actual cusp inertia.

\paragraph{Scope for the original detecting covers.}
An arbitrary cover in F.3.23 has a specific Jacobian
$J_Y$ and the actual CM quotient supplied by F.3.18.
The argument above applies directly when that entire
Jacobian is absolutely simple of dimension three.
It does not replace a proper CM factor by the whole
homology. In fact an original cover of $E_0\setminus O$
still has its actual norm quotient to $E_0$, so a
higher-genus covering Jacobian already has an elliptic
quotient. To use a three-dimensional CM Prym factor
there, one must either construct an actual punctured
genus-three target map or calculate the universal
normalizer trace with the remaining complementary
handles still present. Those maps remain explicit
requirements for the general LM6 A.22 gate.
